\documentclass[10pt]{article}

\usepackage[utf8]{inputenc}

\usepackage[T1]{fontenc}

\usepackage{amsmath}

\usepackage{amsfonts}

\usepackage{amssymb}

\usepackage[version=4]{mhchem}

\usepackage{stmaryrd}

\usepackage{bbold}



\begin{document}

при выполнении условий (1) будет совершать так наз. движение по инерции, напр. двигаться поступательно, равномерно и прямолинейно. Если твердое тело не является свободным (см. Связи механические), то условия его равновесия дают те из равенств (1) (или их следствий), к-рые не содержат реакций наложенных связей; остальные равенства дают уравнения для определения неизвестных реакций. Напр., для тела, имеющего неподвижную ось вращения $O z$, условием равновесия будет



\[

\sum m_{z}\left(F_{k}\right)=0

\]



остальные равенства (1) служат для определения реакций подшипников, закрепляющих ось. Если тело закреплено наложенными связями жестко, то все равенства (1) дают уравнения для определения реакций связей.



Согласно принципу отвердевания, равенства (1), не содержашие реакций внешних связей, дают одновременно необходимые (но недостаточные) условия равновесия любой механич. системы и, в частности, деформируемого тела. Необходимые и достаточные условия равновесия любой механич. системы могут быть найдены с помощью возможных перемещений принципа. Для системы, имеющей $s$ степеней свободы, эти условия состоят в равенстве нулю соответствующих обобщенных сил:



\[

Q_{1}=0, Q_{2}=0, \ldots, Q_{s}=0 .

\]



Из состояний равновесия, определяемых условиями (1) или (2), практически реализуются лишь те, к-рые являются устойчивыми (см. Устойчивость равновесия). $\quad$ С. М. Тарг.



PАВНОВЕСИЯ ДИАТРАММА - см. Фазовая диаграмма. PABHOBECHАЯ MEPA - см. Еикость множества.



PABHOBECHЫE COCTOوHИЯ статистических с истем - состояния статистических систем, в которых средние значения любых термодинамических величин (параметров), характеризующих систему, не зависят от времени, то есть могут оставаться неизменными сколь угодно долго.



Статистич. равновесие не означает равновесия в механич. смысле, то есть и в Р. с. статистич. систем частицы движутся в соответствии с уравнениями Гамильтона, а не покоятся в положения х равновесия (минимумах потенциальной энергии). С этим связаны постоянно существующие небольшие отклонения значений термодинамич. величин (параметров) от равновесных - флуктуации.



В статистич. механике с остоя н ием классич. статистич. системы называется распределение вероятностей (вероятностная мера), заданное на фазовом пространстве системы. Для квантовых систем состояние определяется заданием плотности матрицы с соответствующими условиями симметрии относительно перестановок частиц (см. ФермиДирака статистика, Бозе - Эйнитейна статистика). Эквивалентным способом определения состояния статистич. систем служит задание последовательности распределения вероятностей частичных функций (приведенных матриц плотности - в квантовом случае). В статистич. механике Р.с. статистич. систем описываются с помощью Гиббса pacnределений [микроканонического - для изолированной системы, канонического - для системы, обменивающейся с окружающей средой (термостатом) энергией, и большого канонич.- для системы, обменивающейся с термостатом энергией и частицами].



Используется также понятие локального равновесного состояния статистич. системы (см. Локально возмущенные равновесные функции распределения).



Одной из важных и до конца не решенных проблем статистич. механики является математически строгое обоснование стремления статистич. системы к равновесию, то есть доказательство того, что неравновесные состояния статистич. систем стремятся к Р.с. с течением времени.



Наиболее точное описание термодинамич. свойств статистич. систем в состоянии равновесия может быть получено при исследовании бесконечных систем, к-рые получаются из систем конечного числа частиц в термодинамическом пределе. Их состояния описываются предельными распределениями Гиббса, удовлетворяющими ДЛР-уравнениям [или Майера уравнению, Кирквуда - Зальцбурга уравнению для частичных (приведенных) функций распределения вероятностей].



См. также Термодинамический предел для равновесных состояний непрерывных систем квантовой статистической механики, Термодинамический предел для равновесных состояний непрерывных систем классической статистической механики.



Лит. см. при ст. Гиббсовские состояния. PABHOBECHЫЕ СОСТОЯНИЯ; к и н е т и е с к и й п р е дел - предельный переход по параметрам, входящим в гамильтониан системы частиц, при котором равновесные $n$-частичные функции распределения в пределе факторизуются на одночастичные. Основными примерами кинетич. пределов являются среднего поля предел, Больцмана - Грэда предел (малой плотности), слабой связи и т. п.



Кинетич. предел Р. с. является стационарным (равновесным) решением соответствующего этому пределу кинетического уравнения.



Предел среднего поля состоит в следующем. Пусть



\[

\rho_{\varepsilon}\left(q_{1}, \ldots, q_{m}\right)=\varepsilon^{m} \rho\left(q_{1}, \ldots, q_{m} ; \varepsilon\right), q_{i} \in \mathbb{R}^{v},

\]



a $\rho\left(q_{1}, \ldots, q_{m} ; \varepsilon\right)$ - гиббсовские функции распределения с парным потенциалом $\varepsilon \Phi(q)$, активностью $\varepsilon^{-1} z$, обратной температурой $\beta$. Последовательность $\rho_{\varepsilon}=\left\{\rho_{\varepsilon}\left(q_{1}, \ldots, q_{m}\right)\right\}_{m \geqslant 1}$ удовлетворяет уравнениям Кирквуда - Зальцбурга



где $\alpha=(1,0, \ldots)$



\[

\rho_{\varepsilon}=z\left(K_{\varepsilon} \rho_{\varepsilon}+\alpha\right)

\]



\[

\begin{gathered}

\left(K_{\varepsilon} \rho\right)\left(q_{1}, \ldots, q_{m}\right)= \\

=\exp \left\{-\beta \sum_{k=2}^{m} \Phi\left(q_{1}-q_{k}\right)\right\}\left[\left(1-\delta_{m, 1}\right) \rho\left(q_{2}, \ldots, q_{m}\right)+\right. \\

+\sum_{n \geqslant 1} \frac{1}{n !} \int_{\mathbb{R}^{v n}} d q_{1}^{\prime} \ldots d q_{n}^{\prime} \times \\

\left.\times \prod_{j=1}^{n} f_{\varepsilon}\left(q_{1}-q_{j}^{\prime}\right) \rho\left(q_{2}, \ldots, q_{m}, q_{1}^{\prime}, \ldots, q_{n}^{\prime}\right)\right], \\

f_{\varepsilon}(q)=\frac{1}{\varepsilon} \exp \{-\beta \varepsilon \Phi(q)\}-1, \quad f_{0}(q)=-\beta \Phi(q) .

\end{gathered}

\]



В пределе $\varepsilon \rightarrow 0$ имеет место уравнение



\[

\rho_{0}=z\left(K_{0} \rho_{0}+\alpha\right) \text {. }

\]



Последовательность $\rho_{0}=\left\{\rho_{0}\left(q_{1}, \ldots, q_{m}\right)\right\}_{m \geqslant 1}$ является решением этого уравнения, если



\[

\rho_{0}\left(q_{1}, \ldots, q_{m}\right)=\prod_{i=1}^{m} \rho_{0}\left(q_{i}\right)

\]



и $\rho_{0}(q)$ удовлетворяет уравнению Кирквуда Мон ро



\[

\rho_{0}(q)=z \exp \left\{-\beta \int_{\mathbb{R}^{v}} d q^{\prime} \Phi\left(q-q^{\prime}\right) \rho_{0}\left(q^{\prime}\right)\right\} .

\]



Уравнение Кирквуда - Зальцбурга имеет единственное в пространстве $\mathbb{E}_{\xi}$ решение, если



где



\[

|z| \leqslant \xi^{-1} \exp \left\{-2 \beta B-\xi C_{0}(\beta)\right\}

\]



\[

C_{0}(\beta)=e^{\beta \Phi_{0}} \int_{\mathbb{R}^{v}} d q \Phi(q),

\]



$\Phi(q) \geqslant-\Phi_{0}, B-$ константа из условия устойчивости.



Если потенциал $\Phi(q)$ положителен, то $\rho_{\mathcal{E}}-\rho_{0}$ стремится к нулю в норме пространства $\mathbb{E}_{\xi}$ при $\varepsilon \rightarrow 0$, если $\rho_{\varepsilon}$ и $\rho_{0}-$ соответственно решения допредельного и предельного уравнений Кирквуда - Зальцбурга.



Предел Больцмана - Грэда для бесконечной системы частиц, к-рые взаимодействуют как упругие шары, состоит в том, что диаметр частиц $d \rightarrow 0$, активность $z \rightarrow \infty$,





\end{document}